% ===========================================================================
% VII. Conclusion
% ===========================================================================
\section{Conclusion}
\label{sec:conclusion}

We presented a MARL-based adversarial evaluation framework for a
blockchain Bayesian reputation system designed for additive manufacturing
supply chains.  A 12-action PettingZoo AEC environment models 9 attack
vectors derived from the companion blockchain evaluation, trained with
MAPPO across 11 experimental configurations and 5 seeds each.

The evaluation yields three principal findings.  First, \emph{deterministic
defenses are effective}: self-rating, admin escalation, and gate bypass
attacks are completely deterred (100\% block rate) once adversarial agents
learn through training that the negative expected return is unconditional.
Second, \emph{probabilistic defenses create partial equilibria}: evidence
tampering achieves an $81.5\%$ empirical block rate matching the designed
$p=0.80$ detection probability; the $20\%$ evasion probability prevents
full deterrence but maintains $99.7\%$ honest population behavior.  Third,
\emph{economic deterrence constrains but does not eliminate Sybil and
collusion}: block rates of $52$--$58\%$ limit exploitation without
preventing all adversarial actions, and the honest population converges to
$\geq\!87\%$ honest behavior even under sustained Sybil pressure.

The comprehensive multi-vector attack (C11) identifies the primary open
challenge: with $50\%$ adversarial population and coordinated multi-vector
attacks, the system converges bimodally ($20.0\pm40.0\%$ honest; 4/5 seeds
adversarial-dominant) and the defense system achieves a $48.0\%$ block rate
across 164{,}345 attack attempts.  Addressing this challenge likely
requires extending the observation space with cross-agent comparison signals
and introducing cooperative majority protocols among honest agents.
A terminal reward variant connecting the RL objective to long-term reputation
outcomes eliminates the bimodal failure entirely for C11 ($100.0\%$ honest
across all seeds), identifying outcome-based incentive alignment as the key
design principle for multi-vector defense.

Several directions extend this work.  \emph{Curriculum learning and
adversarial annealing} could address bimodal convergence in C11 by gradually
increasing adversarial population fraction during training, allowing honest
agents to build robust policies before facing full adversarial pressure.
\emph{Enriched observation spaces} that include peer behavior distributions,
network-level reputation entropy, and clustering indicators would enable
agents to detect coordinated Sybil and collusion patterns currently invisible
in the 14-dimensional observation.  \emph{Heterogeneous policy architectures}
that assign separate networks to honest and adversarial agent types would
more faithfully model the asymmetric capabilities of real supply chain actors,
addressing the limitation of MAPPO's shared-policy assumption.
\emph{Population-level coordination mechanisms} --- majority voting signals,
consensus penalties, or group-level rewards --- could provide honest agents
with collective defense capabilities against adversarial clusters, directly
targeting the C11 failure mode.  Finally, \emph{dynamic adversarial
populations} where agents enter and leave the supply chain during training
would test robustness under realistic participation dynamics rather than the
fixed-fraction assumption used here.

This work demonstrates that MARL-based adversarial evaluation is a
productive complement to scripted attack injection for reputation system
security analysis.  The two methodologies agree on deterministic defense
outcomes, while MARL additionally quantifies block rates, reveals honest
convergence dynamics, and discovers emergent multi-vector challenges that
static scripting cannot anticipate.  The complete implementation ---
environment, training harness, configurations, and results --- is released
as an open-source reference for the community.
